\section{Introduction}\label{sec:introduction}

The numerical approximation of high-frequency Helmholtz problems is
challenging because stability and approximation are coupled through the
wavenumber.  Standard low-order finite elements may exhibit a pollution
error that is not explained by best approximation alone, and purely local
residual estimators can substantially underestimate this component before
the asymptotic regime is reached
\cite{BabuskaIhlenburgStrouboulisGangaraj1997a,
BabuskaIhlenburgStrouboulisGangaraj1997b,MelenkSauter2011}.  This has
motivated wavenumber-explicit analyses, enriched and discontinuous
formulations, and adaptive estimators that account for global stability
information; see, among others,
\cite{SauterZech2015,CongreveGedickePerugia2019,
ChaumontFreletErnVohralik2021}.

The localized orthogonal decomposition (LOD) constructs problem-adapted
coarse trial and test functions by correcting standard finite element
functions with fine-scale responses on local patches.  It was introduced
for elliptic multiscale problems in \cite{MalqvistPeterseim2014} and has
since been developed into a flexible numerical homogenization framework
\cite{MalqvistPeterseim2020,EngwerHenningMalqvistPeterseim2019}.  For the
Helmholtz equation, the Petrov--Galerkin LOD eliminates the classical
pollution mechanism under a wavelength-resolution condition and logarithmic
oversampling \cite{Peterseim2017}.  Related developments include two-scale,
multiresolution, and super-localized Helmholtz constructions
\cite{OhlbergerVerfurth2017,HauckPeterseim2022,
FreeseHauckPeterseim2024}.

A posteriori certification of an LOD solution raises a specific difficulty.
The correctors are computed in a finite-dimensional fine space and are then
localized.  An estimator that only tests the residual on that discrete
kernel controls the error to a discrete reference problem, but it does not
see the remaining error between the reference problem and the exact
solution.  Multiplication by a continuous Helmholtz inf--sup constant would
restore reliability, but such a constant is generally unavailable and may
obscure the desired wavenumber robustness.

The present paper avoids both limitations.  Its main contributions are as
follows.

First, we retain a local dual-Riesz estimator on the discrete fine-scale
kernel.  It is locally efficient with respect to the exact error and is
therefore used for inexpensive coarse marking.  Second, we construct a
compatibility-corrected equilibrated flux for the full residual.  The
correction is essential because the LOD solution is not a Galerkin solution
in the fine finite element space.  Third, we introduce total primal and
adjoint corrector defects relative to the continuous kernel.  Their dual
norms are bounded by quotient residual majorants obtained with the same
patchwise mixed technology as the full-residual reconstruction.  Finally,
a computable perturbation argument transfers the inf--sup constant of the
localized Petrov--Galerkin matrix to the continuous ideal LOD spaces.  The
resulting guaranteed upper bound has the form
\[
  \norm{u-U}_\kappa
  \leq
  \left(
    c_W^{-1}
    +
    \frac{\epsA}{(1-\epsA)\betaInf}
  \right)\etaEq.
\]
All quantities on the right are computable.  If the second term is bounded
by a fixed tolerance and the coarse wavelength-resolution condition is
satisfied uniformly, the reliability constant is independent of the
wavenumber.

The exact certificate is more expensive than the local marking indicator.
We therefore evaluate it lazily.  Between certification checkpoints, the
fine mesh and oversampling level are kept fixed and local factorizations are
reused; such an interval is called a cache batch.  This is only an
implementation device and does not introduce a changing reference problem.
When the certificate detects an unresolved corrector defect, a separate
localization diagnostic decides whether to increase the oversampling depth
or to refine the fine mesh.  The analysis establishes guaranteed
reliability and certified termination, but it does not claim estimator
reduction, contraction, or optimal complexity for the complete adaptive
loop.

The paper is organized as follows.  \Cref{sec:model-spaces} introduces the
model problem, the two nested meshes, and the projective quasi-interpolation.
\Cref{sec:lod-framework} develops the continuous and discrete correctors,
the localized Petrov--Galerkin method, and the local marking indicator.
\Cref{sec:exact-certification} constructs the full-residual and total-defect
majorants and proves the exact-solution certificate.
\Cref{sec:adaptive} presents the adaptive algorithm, cache reuse, and cost
model.  The two numerical benchmarks are described in
\cref{sec:numerics}, and \cref{sec:conclusions} concludes the paper.


%======================================================================
